\documentclass[12pt]{article}
\usepackage[margin=1in]{geometry}
\usepackage{enumitem}
\usepackage{titlesec}
\usepackage{parskip}
\usepackage{graphicx}
\usepackage{xcolor}
\usepackage{hyperref}
\usepackage{fontenc}

\title{\textbf{Fearon and Laitin 2003 RG Notes}\\
\large Fearon, James D., and David D. Laitin. ``Ethnicity, Insurgency, and Civil War'' [2003].\\
\textit{American Political Science Review} 97(01): 75-90.}
\author{}
\date{}

\begin{document}

\maketitle

Conventional theories:

(1) high rates of civil wars in 1990s were associated with a shifting international order following the conclusion of the Cold War

(2) greater degrees of diversity (ethnic/religious) alone leads to a higher prevalence of civil war (consider this ``perenialist'', where new and converse theory is ``modernist'' or ``constructivist'' --- that these cleavages become conflict-prone following economic modernization, particularly as upwards mobility becomes both possible and unequally accessible across groups)

(3) future wars can be predicted by evaluating the strongest ongoing ethnic or political grievances

This paper outright disproves (1), qualifies and largely disproves (2) --- disappears once controlling for income --- and does not find strong evidence to support (3)

Biggest factor in civil war outbreak seems to be conditions for insurgency: the efficacy of  small-scale guerilla warfare particularly in rural areas

Financial factors may absolutely be tied to violence, but largely indirectly: economic variables proxy state capacity, such that poorer environments are likely to have less state suppression of violent outbreaks

Core component of insurgency: the rebel group is small (small enough to hide) --- significant outbreaks of violence when this low size is met

Generally suggests that 1990s conflicts are simply the peak of accumulations from the 1950s, particularly as decolonization gave way to Yong (and poor/weak) states which were subsequently ripe for insurgence

Suggests also that t is less important to address/mitigate grievances between ethnicities; they are often simply effects of war, but not necessarily great predictors of whether or not war will arise.

Logical argument (or at least its implications) very similar to Posner 2004: ethnic/religious cleavage's are indeed both present and sometimes salient, but are usually not sufficient on their own to result in widespread violence or political salience. Rather, there must e some additional contributing factors --- usually some degree of inequality or repression enabled by an influx of new wealth or opportunities. In combination between the two papers --- cleavages are the foundations of the causes of violence, rather than the causes of violence themselves

\end{document}
